\section{Contexto, datos y formulación --- Héctor Nieto Colás}

\subsection{Contexto}

Las Jornadas de Ciencia de Datos son un congreso de dos días dividido en ocho
franjas temporales: las franjas 1--4 corresponden al primer día y las 5--8 al
segundo. Se dispone de tres espacios: el Auditorio (200 plazas), el Aula de
ordenadores (40 plazas) y la Sala de seminarios (80 plazas).

La organización dispone de 18 actividades candidatas clasificadas en aprendizaje
automático, estadística e ingeniería de datos. Debe decidir qué actividades se
realizan y en qué espacio y franja comienza cada una. La decisión es real porque
no pueden programarse todas las actividades: el coste total asciende a 12.050
unidades, mientras que el presupuesto disponible es de 6.500.

El contexto aporta las siguientes condiciones:
\begin{itemize}
\item la conferencia inaugural y la de clausura son obligatorias;
\item los talleres y el \emph{hackathon} duran varias franjas y solo pueden
      celebrarse en el Aula de ordenadores;
\item las actividades con asistencia prevista superior a 80 personas no caben
      en la Sala de seminarios;
\item algunos ponentes invitados solo están disponibles en determinadas franjas;
\item existen actividades que compiten por el mismo público o ponente y no deben
      coincidir.
\end{itemize}

\subsection{Datos}

Los datos están separados del código en siete archivos CSV dentro de
\texttt{data/}, de forma que pueda resolverse cualquier otra instancia con el
mismo formato. Las estimaciones de costes, asistencias y valoraciones son los
valores definidos por el grupo para esta instancia.

\begin{table}[htbp]
\centering
\begin{tabular}{p{0.32\textwidth}p{0.58\textwidth}}
\toprule
Archivo & Contenido \\
\midrule
\texttt{actividades.csv} & Duración, categoría, asistencia, valoración,
coste y obligatoriedad \\
\texttt{espacios.csv} & Capacidad de cada espacio \\
\texttt{compat\_espacios.csv} & Espacios compatibles con cada actividad \\
\texttt{disp\_actividades.csv} & Franjas disponibles de cada actividad \\
\texttt{disp\_espacios.csv} & Franjas disponibles de cada espacio \\
\texttt{incompatibilidades.csv} & Pares de actividades incompatibles \\
\texttt{parametros.csv} & Presupuesto, pesos y horizonte temporal \\
\bottomrule
\end{tabular}
\caption{Archivos de datos de la instancia.}
\end{table}

El problema tiene una decisión real porque el coste total supera el presupuesto
y las duraciones suman 27 franjas, más que las 21 ocupaciones disponibles entre
los tres espacios en el horizonte considerado. La instancia contiene dos
actividades obligatorias, varias actividades de varias franjas, disponibilidad
temporal limitada y relaciones de incompatibilidad.

\subsection{Formulación}
\subsubsection{Conjuntos}
\begin{table}[H]
\centering
\begin{tabular}{p{0.20\textwidth}p{0.28\textwidth}p{0.42\textwidth}}
\toprule
\textbf{Símbolo} & \textbf{Definición} & \textbf{Significado} \\
\midrule
$I$ & $\{1,\ldots,18\}$ & Actividades candidatas. \\
$E$ & --- & Espacios disponibles. \\
$T$ & $\{1,\ldots,F\}$ & Franjas temporales. \\
$O$ & $O\subseteq I$ & Actividades obligatorias. \\
$P$ & --- & Pares $(i,j)$ de actividades incompatibles. \\
\bottomrule
\end{tabular}
\caption{Conjuntos del modelo.}
\end{table}

\subsubsection{Parámetros}
\begin{table}[H]
\centering
\begin{tabular}{p{0.20\textwidth}p{0.28\textwidth}p{0.42\textwidth}}
\toprule
\textbf{Símbolo} & \textbf{Definición} & \textbf{Significado} \\
\midrule
$d_i$ & --- & Duración de la actividad $i$, en franjas. \\
$a_i$ & --- & Asistencia esperada de la actividad $i$. \\
$v_i$ & --- & Valoración de la actividad $i$. \\
$c_i$ & --- & Coste de la actividad $i$. \\
$\mathrm{cap}_e$ & --- & Capacidad del espacio $e$. \\
$B$ & --- & Presupuesto máximo. \\
$\mathrm{comp}(i,e)$ & $\{0,1\}$ & Vale 1 si $i$ es compatible con $e$. \\
$\mathrm{dispA}(i,s)$ & $\{0,1\}$ & Vale 1 si $i$ puede celebrarse en $s$. \\
$\mathrm{dispE}(e,s)$ & $\{0,1\}$ & Vale 1 si $e$ está abierto en $s$. \\
$F$ & $8$ & Número de franjas. \\
$FD$ & $4$ & Número de franjas por día. \\
$\mathrm{dia}(t)$ & $\lceil t/FD\rceil$ & Día correspondiente a la franja $t$. \\
$w_A,w_V,w_N$ & $\geq 0$, suma 1 & Pesos de los tres criterios. \\
$N$ & $|I|$ & Número total de actividades candidatas. \\
$A_{\mathrm{tot}}$ & $\sum_{i\in I}a_i$ & Asistencia total de las actividades. \\
$V_{\mathrm{tot}}$ & $\sum_{i\in I}v_i$ & Valoración total de las actividades. \\
\bottomrule
\end{tabular}
\caption{Parámetros del modelo.}
\end{table}

\subsubsection{Variables}
\begin{table}[H]
\centering
\begin{tabular}{p{0.22\textwidth}p{0.20\textwidth}p{0.48\textwidth}}
\toprule
\textbf{Variable} & \textbf{Dominio} & \textbf{Interpretación} \\
\midrule
$x_{i,e,t}$ & $\{0,1\}$ & Vale 1 si la actividad $i$ comienza en el espacio
$e$ en la franja $t$, y 0 en otro caso. \\
$y_i$ & $\{0,1\}$ & Vale 1 si la actividad $i$ se programa, y 0 en otro caso. \\
\bottomrule
\end{tabular}
\caption{Variables binarias del modelo.}
\end{table}

\subsubsection{Función objetivo}
Se maximiza una suma ponderada de asistencia total, valoración total y número
de actividades programadas. La asistencia mide el alcance, la valoración mide
la calidad y el número de actividades evita una programación excesivamente
concentrada en pocas sesiones. Cada componente se normaliza para hacer
comparables sus unidades.

\[
\max Z=
\sum_{i\in I}y_i
\left(
w_A\frac{a_i}{A_{\mathrm{tot}}}
+
w_V\frac{v_i}{V_{\mathrm{tot}}}
+
w_N\frac{1}{N}
\right).
\]

Los pesos son no negativos y cumplen
$w_A+w_V+w_N=1$. Los valores por defecto son $(0.4,0.4,0.2)$; los escenarios
de análisis priorizan, respectivamente, la asistencia $(1,0,0)$, la valoración
$(0,1,0)$ y el número de actividades $(0.1,0.1,0.8)$.

\subsubsection{Restricciones}

\paragraph{R1. Selección única.}
Cada actividad se selecciona exactamente una vez o no se selecciona; por tanto,
no puede aparecer sin asignación ni ocupar varios espacios o franjas de inicio.
\[
y_i=\sum_{e\in E}\sum_{t\in T}x_{i,e,t}
\qquad\forall i\in I.
\]

\paragraph{R2. Actividades obligatorias.}
Las actividades marcadas como obligatorias deben formar parte de la
programación.
\[
y_i=1\qquad\forall i\in O.
\]

\paragraph{R3. Compatibilidad de espacio.}
Una actividad solo puede comenzar en espacios técnicamente compatibles con ella.
\[
x_{i,e,t}=0\qquad
\text{si }\mathrm{comp}(i,e)=0,
\qquad\forall i\in I,\ e\in E,\ t\in T.
\]

\paragraph{R4. Capacidad suficiente.}
La asistencia prevista no puede superar la capacidad del espacio asignado.
\[
x_{i,e,t}=0\qquad
\text{si }a_i>\mathrm{cap}_e,
\qquad\forall i\in I,\ e\in E,\ t\in T.
\]

\paragraph{R5. Disponibilidad de actividad y espacio.}
La disponibilidad se exige para todas las franjas ocupadas:
\[
x_{i,e,t}\leq
\mathrm{dispA}(i,s)\,\mathrm{dispE}(e,s),
\qquad
\forall i\in I,\ e\in E,\ t\in T,
s\in\{t,\ldots,t+d_i-1\}.
\]

\paragraph{R6. Horizonte y límite diario.}
\[
x_{i,e,t}=0
\quad\text{si}\quad
t+d_i-1>F
\quad\text{o}\quad
\mathrm{dia}(t)\neq\mathrm{dia}(t+d_i-1),
\qquad\forall i\in I,\ e\in E,\ t\in T.
\]

\paragraph{R7. Presupuesto.}
El coste total de las actividades seleccionadas no puede superar el presupuesto
disponible.
\[
\sum_{i\in I}c_i y_i\leq B.
\]

\paragraph{R8. Ocupación de un espacio.}
Una actividad que comienza en $t$ ocupa las franjas $t,\ldots,t+d_i-1$:
\[
\sum_{i\in I}
\sum_{t=\max(1,s-d_i+1)}^{s}
x_{i,e,t}
\leq 1
\qquad\forall e\in E,\;\forall s\in T.
\]

\subsubsection{Ampliación}

Como ampliación se incorporan incompatibilidades temporales mediante R9.
Dos actividades incompatibles pueden programarse, pero nunca simultáneamente,
aunque estén asignadas a espacios distintos. Los pares definidos son
$(3,10)$, $(12,16)$, $(9,17)$, $(7,18)$ y $(8,14)$.

La restricción representa actividades que compiten por público o ponente y
afecta tanto a la selección como a la asignación temporal. Por ejemplo, las
actividades 7 y 18 solo están disponibles en la franja 7 y no pueden
programarse conjuntamente.

\paragraph{R9. Incompatibilidad entre actividades.}
Las parejas incompatibles no pueden desarrollarse simultáneamente, aunque se
encuentren en espacios diferentes, para evitar competir por público o ponente.
Para cada par $(i,j)\in P$ y para cada franja $s\in T$:
\[
\sum_{e\in E}
\sum_{t=\max(1,s-d_i+1)}^{s}x_{i,e,t}
+
\sum_{e\in E}
\sum_{t=\max(1,s-d_j+1)}^{s}x_{j,e,t}
\leq 1.
\]

